% !TEX program = xelatex
% 공통수학2 · Ⅰ. 도형의 방정식 · 2. 원의 방정식 · 02 원과 직선의 위치 관계 · 2/2차시
% 학생용 활동지 (답안 없음)
\documentclass[11pt,a4paper]{article}
\usepackage{kotex}
\usepackage{iftex}
\ifPDFTeX\else
  \setmainhangulfont{Noto Serif CJK KR}
  \setsanshangulfont{Noto Sans CJK KR}
\fi
\usepackage[margin=2.1cm]{geometry}
\usepackage{parskip}
\usepackage{enumitem}
\usepackage{array}
\usepackage{tabularx}
\usepackage{xcolor}
\usepackage{amssymb}
\usepackage{tikz}
\usepackage[most]{tcolorbox}
\usepackage{fancyhdr}
\usepackage{titlesec}

\definecolor{goldc}{HTML}{B07C22}
\definecolor{goldstrong}{HTML}{8A5F16}
\definecolor{indigoc}{HTML}{2B3B57}
\definecolor{indigosoft}{HTML}{6E82A6}
\definecolor{linec}{HTML}{D6DACB}
\definecolor{surface2}{HTML}{E4E8DC}
\definecolor{writeline}{HTML}{C7CCBA}
\definecolor{inksoft}{HTML}{52604F}

\titleformat{\section}{\normalfont\sffamily\bfseries\large\color{indigoc}}{\thesection}{0.6em}{}
\titlespacing{\section}{0pt}{14pt}{6pt}
\newtcolorbox{goalbox}{enhanced, breakable, colback=white, colframe=goldc, boxrule=0.5pt, leftrule=3pt, arc=2pt, left=10pt, right=10pt, top=8pt, bottom=8pt}
\newtcolorbox{sheetbox}{enhanced, breakable, colback=white, colframe=linec, boxrule=0.6pt, arc=3pt, left=11pt, right=11pt, top=9pt, bottom=9pt}
\newcommand{\writespace}[1]{%
  \par\nobreak\vspace{3pt}
  \foreach \n in {1,...,#1} {%
    \noindent\textcolor{writeline}{\rule{\linewidth}{0.4pt}}\par\vspace{0.62cm}%
  }
}
\newcommand{\fillblank}[1]{\makebox[#1]{\hrulefill}}
\pagestyle{fancy}\fancyhf{}\renewcommand{\headrulewidth}{0pt}
\fancyfoot[L]{\footnotesize\color{inksoft} 공통수학2 · 2026학년도 2학기 · 지평선고등학교}
\fancyfoot[R]{\footnotesize\color{inksoft} 다음 시간 --- 중단원 마무리(원의 방정식)}

\begin{document}

\noindent
\begin{minipage}[b]{0.62\linewidth}
  {\sffamily\footnotesize\color{indigosoft} 공통수학2 · Ⅰ. 도형의 방정식 · 2. 원의 방정식}\\[2pt]
  {\sffamily\bfseries\Large 02 원과 직선의 위치 관계 \textperiodcentered\ 12차시}
\end{minipage}\hfill
\begin{minipage}[b]{0.36\linewidth}
  \raggedleft\small\color{inksoft}
  반\ \fillblank{1.6cm} \quad 번호\ \fillblank{1.6cm}\\[4pt]
  이름\ \fillblank{3.6cm}
\end{minipage}

\vspace{4pt}
\noindent{\color{goldc}\rule{\linewidth}{2.2pt}}
\vspace{10pt}

\begin{goalbox}
\textbf{오늘의 목표}\\
기울기가 주어진 경우, 원 위의 점, 원 밖의 점 --- 세 가지 상황에서 원의 접선의 방정식을 구할 수 있다.
\vspace{4pt}
\small\color{inksoft}
$\square$ 자신 있음 \quad $\square$ 조금 헷갈림 \quad $\square$ 처음 봄 \hfill (수업 전 나의 상태에 표시)
\end{goalbox}

\section{개념 정리 --- 세 가지 접선}

\begin{sheetbox}
\textbf{1) 기울기가 주어질 때} --- 원 $x^2+y^2=r^2$에 접하고 기울기가 $m$인 직선의 방정식은
\[
y = mx \pm \fillblank{3.2cm}
\]

\vspace{6pt}
\textbf{2) 원 위의 점에서} --- 원 $x^2+y^2=r^2$ 위의 점 P$(x_1,y_1)$에서의 접선의 방정식은
\[
\fillblank{2cm}\,x + \fillblank{2cm}\,y = r^2
\]

\vspace{6pt}
\textbf{3) 원 밖의 점에서} --- 접점을 P$(x_1,y_1)$이라 놓고, (①) P에서의 접선이 주어진 점을 지난다는 조건과 (②) P가 원 위의 점이라는 조건, 두 식을 \fillblank{3.2cm}하여 $x_1,y_1$을 구한다.
\end{sheetbox}

\section{함께 채우기 --- 원 밖의 점에서 그은 접선}

\begin{sheetbox}
점 $(4,0)$에서 원 $x^2+y^2=4$에 그은 접선의 방정식을 구해 보자. 접점을 P$(x_1,y_1)$이라 하면 접선의 방정식은
\[
x_1x + y_1y = 4
\]
이 직선이 점 $(4,0)$을 지나므로 $x_1 = \fillblank{1.6cm}$

또, 점 P가 원 위의 점이므로 $x_1^2+y_1^2=4$ \; ...①

$x_1=\fillblank{1.6cm}$을 ①에 대입하면 $y_1 = \pm\fillblank{1.6cm}$

따라서 구하는 접선의 방정식은
\[
x + \fillblank{1.6cm}\,y = 4\,, \qquad x - \fillblank{1.6cm}\,y = 4
\]
\end{sheetbox}

\section{연습 문제}

\begin{sheetbox}
\textbf{1.} 다음을 구하시오.\\
\hspace*{1.6em}⑴ 원 $x^2+y^2=2$에 접하고 기울기가 $-1$인 직선의 방정식\\
\hspace*{1.6em}⑵ 원 $x^2+y^2=9$ 위의 점 $(0,-3)$에서의 접선의 방정식
\writespace{3}

\vspace{6pt}
\textbf{2.} 점 $(3,-1)$에서 원 $x^2+y^2=5$에 그은 접선의 방정식을 구하시오.
\writespace{5}
\end{sheetbox}

\section{사고의 가시화 --- Connect · Extend · Challenge}

\noindent{\footnotesize\color{inksoft} `원과 직선의 위치 관계' 전체(11$\sim$12차시)를 떠올리며 아래 세 칸을 채워 보자.}\\[6pt]
\begin{tabularx}{\linewidth}{@{}XXX@{}}
\textbf{\footnotesize\color{indigoc} Connect --- 무엇과 연결되나?} &
\textbf{\footnotesize\color{indigoc} Extend --- 어디까지 확장됐나?} &
\textbf{\footnotesize\color{indigoc} Challenge --- 아직 궁금한 것은?} \\[4pt]
\end{tabularx}
\noindent
\begin{minipage}[t]{0.32\linewidth}\writespace{3}\end{minipage}\hfill
\begin{minipage}[t]{0.32\linewidth}\writespace{3}\end{minipage}\hfill
\begin{minipage}[t]{0.32\linewidth}\writespace{3}\end{minipage}

\section{마무리 성찰}

\begin{sheetbox}
\begin{minipage}[t]{0.48\linewidth}
{\footnotesize\color{inksoft} `원과 직선의 위치 관계' 소단원을 마치며 한 문장으로}
\writespace{2}
\end{minipage}\hfill
\begin{minipage}[t]{0.48\linewidth}
{\footnotesize\color{inksoft} 감사 일기 / 유념 한 줄}
\writespace{2}
\end{minipage}
\end{sheetbox}

\end{document}
